% =============================================================================
% notacion.tex · Notación única de la tesis (macros y lista de símbolos)
% Regla: antes de introducir un símbolo se comprueba que no exista aquí. Si se cambia
% uno, se actualiza en toda la tesis y en el informe (tesis/), que deriva de esta lista.
% =============================================================================

% ---- Conjuntos y operadores --------------------------------------------------
\newcommand{\RR}{\mathbb{R}}                       % números reales
\newcommand{\NN}{\mathbb{N}}                       % números naturales
\newcommand{\T}{\top}                              % transpuesta
\DeclareMathOperator{\tr}{tr}                      % traza
\DeclareMathOperator{\diag}{diag}                  % matriz diagonal / diagonal de una matriz
\DeclareMathOperator{\rango}{rango}                % rango de una matriz
\DeclareMathOperator*{\argmin}{arg\,min}
\DeclareMathOperator*{\argmax}{arg\,max}
\DeclareMathOperator{\RV}{RV}                      % coeficiente RV
\DeclareMathOperator{\CH}{CH}                      % índice de Caliński y Harabasz
\DeclareMathOperator{\KMO}{KMO}                    % medida de Kaiser-Meyer-Olkin
\DeclareMathOperator{\congr}{cong}                 % congruencia de Tucker
\DeclareMathOperator{\ctr}{ctr}                    % contribución a un eje
\DeclareMathOperator{\qlt}{cal}                    % calidad de representación en un plano
\newcommand{\norm}[1]{\lVert #1\rVert}             % norma euclídea
\newcommand{\normF}[1]{\lVert #1\rVert_{F}}        % norma de Frobenius
\newcommand{\prodF}[2]{\langle #1,#2\rangle_{F}}   % producto interno de Frobenius
\newcommand{\abs}[1]{\lvert #1\rvert}
\newcommand{\prom}[1]{\langle #1\rangle}           % media temporal de una sucesión en un tramo
\DeclareMathOperator{\vart}{var}                   % varianza temporal de una sucesión en un tramo

% ---- Datos -------------------------------------------------------------------
\newcommand{\mX}{\mathbf{X}}       % matriz de datos observados, n x p
\newcommand{\mZ}{\mathbf{Z}}       % matriz de datos centrados y estandarizados, n x p
\newcommand{\vx}{\mathbf{x}}       % fila de X (vector del mes t)
\newcommand{\vz}{\mathbf{z}}       % fila o columna de Z
\newcommand{\mI}{\mathbf{I}}       % identidad
\newcommand{\vuno}{\mathbf{1}}     % vector de unos

% ---- PCA ----------------------------------------------------------------------
\newcommand{\mR}{\mathbf{R}}           % matriz de correlaciones
\newcommand{\mV}{\mathbf{V}}           % vectores propios de R
\newcommand{\vv}{\mathbf{v}}           % vector propio
\newcommand{\mLambda}{\boldsymbol{\Lambda}} % valores propios en diagonal
\newcommand{\mE}{\mathbf{E}}           % matriz p x d con columnas ortonormales (candidata)
\newcommand{\mP}{\mathbf{P}}           % proyector ortogonal
\newcommand{\vc}{\mathbf{c}}           % componente principal (scores)
\newcommand{\mM}{\mathbf{M}}           % aplicación lineal de reconstrucción (rango <= d)
\newcommand{\mO}{\mathbf{O}}           % matriz ortogonal
\newcommand{\mN}{\mathbf{N}}           % matriz genérica (teorema de Eckart-Young)

% ---- EPCA ---------------------------------------------------------------------
\newcommand{\nw}{n_{w}}                % ancho de la ventana
\newcommand{\Wt}{\mathcal{W}}          % conjunto de meses de una ventana
\newcommand{\vm}{\mathbf{m}}           % media ponderada de la ventana
\newcommand{\mC}{\mathbf{C}}           % covarianza ponderada de la ventana
\newcommand{\mD}{\mathbf{D}}           % diagonal de C

% ---- HFC ----------------------------------------------------------------------
\newcommand{\gA}{\mathcal{A}}          % grupo de variables
\newcommand{\gB}{\mathcal{B}}          % grupo de variables
\newcommand{\fac}{\boldsymbol{\varphi}}% factor (representante) de un grupo
\newcommand{\dif}{\boldsymbol{\psi}}   % componente de diferencias de un nodo

% ---- Partición y comparación -------------------------------------------------
\newcommand{\vy}{\mathbf{y}}           % punto de una trayectoria a particionar
\newcommand{\mF}{\mathbf{F}}           % configuración (RV, Procrustes)
\newcommand{\estr}{\mathcal{S}}        % estructura estimada de un conjunto de meses
\newcommand{\interv}{\mathcal{I}}      % intervalo de un régimen

% ---- Episodios de referencia y validación externa (iteración 3b) --------------
\newcommand{\epi}{\mathsf{e}}          % episodio de referencia
\newcommand{\Epis}{\mathcal{E}}        % conjunto de episodios de referencia
\newcommand{\tini}[1]{\underline{t}_{#1}} % mes de inicio de un episodio
\newcommand{\tfin}[1]{\overline{t}_{#1}}  % mes de fin (efectivo) de un episodio

% ---- Ventanas, nulos y criterios de la parte dinámica (iteración 3b) ---------
\newcommand{\nomega}{n_{\omega}}       % n efectivo de Kish de los pesos de una ventana
\newcommand{\nhorn}{n_{\mathrm{H}}}    % filas del análisis paralelo de Horn por ventana
\newcommand{\Kest}{K^{\ast}}           % número de segmentos (clases) elegido por un criterio
\newcommand{\Kef}{K_{\mathrm{ef}}}     % mayor número de segmentos factible con n_min
\newcommand{\Lb}{L_{\mathrm{b}}}       % longitud de bloque del remuestreo (no confundir con L del LSTM)
\newcommand{\mb}{m_{\mathrm{b}}}       % número de bloques de una réplica
\newcommand{\RVad}[1]{\RV^{\mathrm{ad}}_{#1}}            % RV entre ventanas adyacentes no solapadas
\newcommand{\umbralRV}[1]{\underline{\RV}^{\,\mathrm{ad}}_{#1}} % p5 por posición del nulo
\newcommand{\RVadr}[2]{\RV^{\mathrm{ad},(#2)}_{#1}}       % RV adyacente de la réplica #2
\newcommand{\umbralRVr}[2]{\underline{\RV}^{\,\mathrm{ad},(#2)}_{#1}} % p5 sin la réplica
\newcommand{\ind}{\mathbb{I}}          % función indicadora
\newcommand{\pib}{\pi^{\mathrm{b}}}    % sucesión de índices por bloques (no confundir con la constante pi)
\DeclareMathOperator{\fr}{fr}          % frontera de un par de ventanas adyacentes
\DeclareMathOperator{\red}{red}        % reducción relativa del costo de Fisher
\newcommand{\pval}{\mathrm{pv}}        % p-valor de Monte Carlo
\newcommand{\paso}{\mathsf{s}}         % paso entre ventanas de la regla del 06

% ---- Aprendizaje profundo ----------------------------------------------------
\newcommand{\mA}{\mathbf{A}}           % codificador lineal
\newcommand{\mB}{\mathbf{B}}           % decodificador lineal
\newcommand{\veta}{\boldsymbol{\eta}}  % código latente
\newcommand{\vh}{\mathbf{h}}           % estado oculto del LSTM
\newcommand{\vcel}{\mathbf{s}}         % estado de la celda de memoria del LSTM
\newcommand{\vig}{\mathbf{g}^{\mathrm{e}}} % compuerta de entrada
\newcommand{\vfg}{\mathbf{g}^{\mathrm{o}}} % compuerta de olvido
\newcommand{\vog}{\mathbf{g}^{\mathrm{s}}} % compuerta de salida
\newcommand{\mW}{\mathbf{W}}           % pesos de entrada del LSTM
\newcommand{\mU}{\mathbf{U}}           % pesos recurrentes del LSTM
\newcommand{\vb}{\mathbf{b}}           % sesgo

% =============================================================================
% Lista de símbolos (se imprime en las páginas preliminares con \listadesimbolos)
% =============================================================================
\newcommand{\gruposimbolos}[1]{\multicolumn{2}{@{}l}{\textbf{#1}}\\[2pt]}
\newcommand{\listadesimbolos}{%
\chapter*{Lista de símbolos}
\addcontentsline{toc}{chapter}{Lista de símbolos}
\markboth{Lista de símbolos}{Lista de símbolos}
{\small
\begin{longtable}{@{}p{0.24\textwidth}p{0.70\textwidth}@{}}
\toprule
Símbolo & Significado \\
\midrule
\endhead
\bottomrule
\endfoot
\gruposimbolos{Generales}
$\RR$, $\NN$, $\mathbb{Z}$ & Números reales, naturales y enteros. \\
$\mI_p$, $\vuno$ & Matriz identidad de orden $p$; vector de unos. \\
$\mathbf{a}^{\T}$, $\tr$, $\diag$, $\rango$ & Transpuesta, traza, matriz diagonal (o diagonal de una matriz) y rango. \\
$\norm{\cdot}$, $\normF{\cdot}$, $\prodF{\cdot}{\cdot}$ & Norma euclídea; norma y producto interno de Frobenius, $\prodF{\mathbf{G}}{\mathbf{H}}=\tr(\mathbf{G}^{\T}\mathbf{H})$. \\
$\lfloor\cdot\rfloor$, $\lceil\cdot\rceil$ & Parte entera inferior y superior. \\
$\abs{\cdot}$, $\operatorname{sgn}$ & Valor absoluto o módulo de un número complejo; función signo. \\
$\mathrm{i}$, $e$ & Unidad imaginaria; base de los logaritmos naturales ($\ln$). \\
$O(\cdot)$ & Orden del número de operaciones de un algoritmo. \\
$\ind\{\cdot\}$, $\Pr$, $\mathbb{E}$, $\#$ & Función indicadora; probabilidad; esperanza; número de elementos de un conjunto. \\
$\mathbf{a}$, $\mathbf{b}$, $\xi$ & Vectores y escalares genéricos en enunciados y demostraciones. \\
$\mathbf{G}\succeq 0$ & Matriz simétrica semidefinida positiva. \\
$\odot$, $\sigma(\cdot)$, $\tanh$ & Producto elemento a elemento; función logística; tangente hiperbólica. \\[4pt]
\gruposimbolos{Datos}
$n$, $p$ & Número de meses (observaciones) y de series (variables); en el análisis oficial, $n=\nMesesNiveles$ y $p=\nSeries$ (\cref{sec:met-datos}). \\
$t$, $s$ & Índices de mes, $t,s\in\{1,\dots,n\}$. \\
$j$, $j'$ & Índices de serie, $j,j'\in\{1,\dots,p\}$. \\
$\mX$, $\vx_t$, $x_{tj}$ & Matriz de datos $n\times p$ (series sin transformar o transformadas); su fila $t$; su elemento. \\
$\mZ$, $\vz_t$, $\vz_j$, $z_{tj}$ & Datos centrados y estandarizados con divisor $n$; fila $t$; columna $j$; elemento. \\
$\bar x_j$, $\hat\sigma_j$ & Media y desviación típica (divisor $n$) de la serie $j$. \\
$\Delta$ & Primera diferencia, $\Delta x_t=x_t-x_{t-1}$; para vectores, $\Delta\vx_t=\vx_t-\vx_{t-1}$. \\
$\hat x_t$, $\hat\epsilon_t$ & Recta de mínimos cuadrados de una serie sobre el tiempo y sus residuos (regla D5, nulo mixto). \\
$P$, $\vartheta$, $\gamma(P)$ & Período de una oscilación en meses; frecuencia angular $2\pi/P$; ganancia de $\Delta$. \\
$\varrho$, $\psi_0$ & Amplitud y fase de una oscilación. \\
$t_0$, $T$ & Mes anterior al primero de un tramo de meses consecutivos y longitud del tramo (\cref{prop:ganancia}). \\
$\prom{y}_{t_0,T}$, $\vart_{t_0,T}(y)$ & Media y varianza temporales de la sucesión $(y_t)$ en los meses $t_0+1,\dots,t_0+T$. \\
$y$, $\varrho_y$, $\psi_y$ & Oscilación genérica, su amplitud y su fase, en la demostración de la \cref{prop:ganancia}. \\[4pt]
\gruposimbolos{Análisis de componentes principales}
$\mR$, $r_{jj'}$ & Matriz de correlaciones $\mZ^{\T}\mZ/n$ y su elemento. \\
$\lambda_\ell$, $\vv_\ell$, $v_{j\ell}$ & Valor propio $\ell$ de $\mR$ (en orden decreciente), vector propio unitario asociado y su coordenada $j$. \\
$\ell$, $d$ & Índice de componente; número de componentes retenidas. \\
$\mV$, $\mLambda$, $\mV_d$ & Matriz de vectores propios; matriz diagonal de valores propios; primeras $d$ columnas de $\mV$. \\
$\mE$, $\mP$ & Matriz $p\times d$ con columnas ortonormales; proyector ortogonal $\mE\mE^{\T}$. \\
$J(\mE)$ & Error cuadrático medio de reconstrucción del PCA. \\
$w_\ell$ & En la demostración de la \cref{prop:pca}, $\vv_\ell^{\T}\mP\vv_\ell$. \\
$\vc_\ell$, $c_{t\ell}$ & Componente principal $\ell$ ($\mZ\vv_\ell$) y su valor en el mes $t$. \\
$\rho_{j\ell}$ & Correlación entre la serie $j$ y la componente $\ell$. \\
$\cos^2_{j\ell}$, $\ctr_{j\ell}$, $\ctr_{t\ell}$ & Calidad de representación de la serie $j$ en el eje $\ell$; contribución de la serie $j$ y del mes $t$ al eje $\ell$. \\
$\qlt_j(\ell,\ell')$ & Calidad de representación de la serie $j$ en el plano $(\ell,\ell')$. \\
$\mN$, $\varsigma_\ell$ & Matriz genérica que se aproxima (teorema de Eckart y Young) y sus valores singulares. \\
$\mM$, $\mO$ & Aplicación lineal de reconstrucción de rango a lo sumo $d$; matriz ortogonal. \\
$\lambda^{\mathrm{H}}_\ell$, $\lambda^{\mathrm{H}}_\ell(n',p)$, $\beta_\ell$ & Umbral del análisis paralelo de Horn, en general o con una tabla de $n'$ filas y $p$ columnas; valor esperado de la vara quebrada. \\
$r^{jj'}$, $q_{jj'}$ & Elemento de $\mR^{-1}$; correlación anti-imagen (parcial) entre las series $j$ y $j'$. \\
$\KMO$, $\KMO_j$ & Medida global e individual de Kaiser-Meyer-Olkin. \\
$\chi^2_{\mathrm{B}}$ & Estadístico de esfericidad de Bartlett. \\[4pt]
\gruposimbolos{Clasificación Factorial Jerárquica}
$\gA$, $\gB$ & Grupos (clases) de variables. \\
$\fac_{\gA}$ & Factor (representante) del grupo $\gA$, vector de $\RR^n$. \\
$\dif_{\gA\cup\gB}$ & Componente de diferencias del nodo que fusiona $\gA$ y $\gB$. \\
$\mu_1^{\gA\gB}\ge\mu_2^{\gA\gB}$ & Valores propios de la matriz de covarianzas de $[\fac_{\gA}\ \fac_{\gB}]$; $\mu_2^{\gA\gB}$ es el índice de fusión. \\
$(\alpha,\beta)$ & Vector propio unitario asociado a $\mu_1^{\gA\gB}$ (coeficientes de la fusión). \\
$\kappa_{\gA}$, $\kappa_{\gB}$, $\kappa_{\gA\gB}$, $\mathbf{K}_{\gA\gB}$ & Varianzas de $\fac_{\gA}$ y $\fac_{\gB}$, su covarianza y la matriz $2\times2$ que forman. \\
$\mu_1^{\text{raíz}}$ & Primer valor propio de la última fusión. \\
$\mR_{\gA}$, $\lambda_1(\mR_{\gA})$, $\mathbf{w}^{\gA}$ & Submatriz de $\mR$ de las series de $\gA$; su mayor valor propio; coeficientes de $\fac_{\gA}$ sobre las $\vz_j$, $j\in\gA$. \\[4pt]
\gruposimbolos{Análisis evolutivo}
$\nw$ & Ancho de la ventana móvil, en meses. \\
$N$ & Número de ventanas, $N=n-\nw+1$; en general, longitud de una trayectoria. \\
$i$, $t_i$ & Índice de la ventana que empieza en el mes $i$; mes al que se asigna, $t_i=i+\lfloor\nw/2\rfloor$. \\
$\Wt_t$ & Conjunto de meses de la ventana asignada al mes $t$. \\
$u$, $\omega_u$ & Posición dentro de la ventana, $u=1,\dots,\nw$, y su peso. \\
$\nu$ & Parámetro de forma de los pesos gaussianos (argumento \texttt{nsd} de \texttt{aede}). \\
$\vm_t$, $\mC_t$, $\mD_t$, $\mR_t$ & Media, covarianza, diagonal de la covarianza y correlación ponderadas de la ventana. \\
$\lambda_\ell(t)$, $\vv_\ell(t)$ & Valor y vector propio $\ell$ de $\mR_t$. \\
$\tilde\omega_u$, $\tilde\vz_s$ & Peso sin normalizar; mes $s$ estandarizado con la media y la varianza ponderadas. \\
$\mathbf{G}_b$ & Matriz diagonal de factores de escala $b_j>0$ (con desplazamientos $a_j$). \\
$\nomega$ & Tamaño efectivo de Kish de los pesos de la ventana, $(\sum_u\omega_u^2)^{-1}$ (\cref{def:kish}). \\
$c(\nu)$, $a_\nu$ & Límite de $\nomega/\nw$ con pesos gaussianos y $a_\nu=\nu/\sqrt2$ (\cref{prop:kish}). \\
$\xi_u$, $S_k$, $I_k$ & Puntos equiespaciados, sumas de Riemann e integrales de la demostración de la \cref{prop:kish}. \\
$\nhorn$ & Filas del análisis paralelo por ventana, $\lceil\nomega\rceil$ (\cref{def:umbral-ventana}). \\
$\paso$, $M$, $\ell_0$ & Paso entre ventanas, número de ventanas y eje que sigue la regla del 06 (\cref{def:regla06}). \\
$\mathbf{a}_k$, $\tilde{\mathbf{a}}_k$ & Vector propio unitario del eje $\ell_0$ en la ventana $k$ de la regla del 06, sin alinear y alineado. \\
$\rho(\cdot,\cdot)$, $\operatorname{sgn}^{+}$ & Correlación de Pearson entre las coordenadas de dos vectores; signo con $\operatorname{sgn}^{+}(0)=1$. \\
$\rho^{(06)}_k$, $\bar\rho$, $s_\rho$ & Estabilidad de la regla del 06 en la ventana $k$; su media y su desviación típica. \\
$c^{\mathrm{ent}}_{\epi}$, $c^{\mathrm{sal}}_{\epi}$ & Centros de entrada y de salida del episodio $\epi$ en la ventana (\cref{def:borde}). \\[4pt]
\gruposimbolos{Partición de Fisher y número de clases}
$\vy_t\in\RR^{q}$, $q$ & Punto $t$ de la trayectoria que se particiona; su dimensión. \\
$K$, $k$ & Número de segmentos (clases); índice de segmento. \\
$\tau_k$ & Último punto del segmento $k$, con $0=\tau_0<\tau_1<\dots<\tau_K=N$. \\
$n_{\min}$ & Longitud mínima de un segmento. \\
$Q(a,b)$, $\bar\vy_{a:b}$ & Suma de cuadrados intrasegmento de los puntos $a,\dots,b$ y su media. \\
$Q^\ast_K(b)$ & Mínimo de la suma de cuadrados intra al dividir los puntos $1,\dots,b$ en $K$ segmentos. \\
$T$ & Suma de cuadrados total de la trayectoria. \\
$K_{\max}$, $Q_K$ & Número máximo de segmentos; $Q^\ast_K(N)$. \\
$\Kef$ & Mayor número de segmentos factible, $\min(K_{\max},\lfloor N/n_{\min}\rfloor)$ (\cref{cor:fisher-factible}). \\
$\Kest$, $\Kest_{\mathrm{codo}}$, $\Kest_{\mathrm K}$ & Número de segmentos que elige un criterio (el $k^\ast$ del código); el del codo por caída máxima y el de la regla K (\cref{def:criterios-k}). \\
$\red_K$, $\upsilon$ & Reducción relativa del costo al pasar de $K$ a $K+1$ segmentos; umbral de la regla K. \\
$w_{\mathrm m}$, $\bar c_{t\ell}$ & Longitud de la media móvil causal y \emph{score} suavizado de la datación PCA-MM12 (\cref{def:datacion}). \\
$\CH^{\mathrm{h}}_K$, $\CH^{\mathrm{a}}_K$ & Índice de Caliński y Harabasz en su forma habitual y en la variante de \texttt{aede}. \\
$K^{\mathrm h}$, $K^{\mathrm a}$ & Número de segmentos elegido por cada variante. \\[4pt]
\gruposimbolos{Régimen, transición y comparación de estructuras}
$\estr(\cdot)$ & Estructura estimada de un conjunto de meses. \\
$\tau$ & Mes de una transición (\cref{def:transicion}); no confundir con $\tau_k$. \\
$\Wt_\tau^{-}$, $\Wt_\tau^{+}$ & Conjuntos de meses contiguos adyacentes, antes y después de $\tau$. \\
$\delta(\cdot,\cdot)$, $\varepsilon$ & Disimilitud entre estructuras; tolerancia. \\
$\interv_k$, $\Wt$, $\Wt'$ & Intervalo de un régimen; conjuntos de meses contiguos (ventanas). \\
$\mF_1$, $\mF_2$, $m$ & Configuraciones no nulas de las mismas $m$ filas (centradas o no, \cref{def:rv}). \\
$\mF_t$, $\mR^{[d]}_t$ & Matriz $p\times d$ de cargas de la ventana de centro $t$ y su producto $\mF_t\mF_t^{\T}$ (\cref{def:rv-cargas}). \\
$\RV^{(d)}(t,t')$, $\rho_{\mR}(t,t')$ & RV de cargas sin centrar con $d$ ejes; coeficiente entre las matrices de correlación completas. \\
$\fr(i)$, $\RVad{i}$ & Frontera y RV del par $i$ de ventanas adyacentes no solapadas (\cref{def:pares}). \\
$\RV(\mF_1,\mF_2)$ & Coeficiente RV. \\
$c$ & Constante de proporcionalidad positiva en la \cref{prop:rv}(a). \\
$a_1$, $a_2$ & Factores de escala no nulos en la \cref{prop:rv}(b). \\
$\mathbf{G}$, $\mathbf{H}$ & En la demostración de la \cref{prop:rv}, $\mF_1\mF_1^{\T}$ y $\mF_2\mF_2^{\T}$. \\
$\mO_1$, $\boldsymbol{\Sigma}$, $\mO_2$, $\boldsymbol{\Theta}$ & Descomposición en valores singulares de $\mF_1^{\T}\mF_2$; matriz ortogonal auxiliar de la \cref{prop:procrustes}. \\
$\mO^\ast$ & Rotación de Procrustes. \\
$\congr(\mathbf{a},\mathbf{b})$ & Coeficiente de congruencia de Tucker entre dos vectores. \\[4pt]
\gruposimbolos{Episodios de referencia y nulos}
$\epi$, $\Epis$ & Episodio de referencia; conjunto de episodios. \\
$\tini{\epi}$, $\tfin{\epi}$ & Mes de inicio y mes de fin (efectivo) del episodio $\epi$ (\cref{def:episodio}). \\
$h$ & Tolerancia en meses (\cref{def:tolerancia}). \\
$\tau^\circ$ & Mes de una transición que se quiere confirmar (\cref{def:confirmada}). \\
$\Lb$, $\mb$, $\iota_k$ & Longitud de bloque, número de bloques e inicios de bloque (\cref{def:bloques}); $\Lb$ no es la $L$ del LSTM. \\
$\pib(t)$, $n'$ & Sucesión de índices por bloques; su longitud. \\
$\dot x_t$, $\overline{\Delta x}_j$ & Incrementos centrados y media de los incrementos de la serie $j$ (\cref{def:hac}). \\
$\hat\Gamma_k$, $m_{\mathrm{HAC}}$, $\bar m$ & Autocovarianza de rezago $k$ de los incrementos; número de rezagos; $m_{\mathrm{HAC}}+1$. \\
$\hat\sigma^2_{\mathrm{HAC},j}$, $t^{\mathrm{HAC}}_j$ & Varianza de largo plazo y estadístico $t$ de la deriva de la serie $j$. \\
$x^\ast_{tj}$ & Serie simulada del nulo de tendencias mixto (\cref{def:nulo-mixto}). \\
$B$, $\zeta$, $\mX^{(\zeta)}$ & Número de réplicas; índice de réplica; réplica del nulo de estructura (\cref{def:nulo-estructura}). \\
$\RVadr{i}{\zeta}$, $\umbralRV{i}$, $\umbralRVr{i}{-\zeta}$ & RV adyacente de la réplica $\zeta$; cuantil de 5\,\% por posición; el mismo sin la réplica $\zeta$ (\cref{def:oe4}). \\
$\varpi_0$, $\varpi_\zeta$, $\pval$ & Proporción de pares bajo el umbral, observada y de la réplica $\zeta$; p-valor del criterio de OE4. \\
$\mathcal{P}_j$, $\tilde\varpi_j$, $\widetilde{\pval}$, $\mathcal{J}$, $\alpha_0$ & Perfiles, proporciones y p-valor de la versión simétrica; conjunto y nivel auxiliares (\cref{prop:oe4}). \\
$\mathcal{T}_a$, $E_t$, $E^{(95)}_a$ & Meses de entrenamiento del año $a$; error de reconstrucción del mes $t$; su percentil 95 en $\mathcal{T}_a$ (\cref{def:walkforward}). \\[4pt]
\gruposimbolos{Aprendizaje profundo}
$f_\theta$, $g_\phi$ & Codificador y decodificador, con parámetros $\theta$ y $\phi$. \\
$\veta_t\in\RR^{d}$ & Código latente del mes $t$. \\
$\Omega(\theta,\phi)$ & Término de regularización. \\
$\mA$, $\mB$ & Codificador lineal $d\times p$ y decodificador lineal $p\times d$. \\
$L$ & Longitud de la ventana que lee el LSTM. \\
$\vh_t\in\RR^{n_h}$, $n_h$ & Estado oculto del LSTM y su dimensión. \\
$\vcel_t$, $\tilde\vcel_t$ & Estado de la celda de memoria y candidato. \\
$\vig_t$, $\vfg_t$, $\vog_t$ & Compuertas de entrada, de olvido y de salida. \\
$\mW_\bullet$, $\mU_\bullet$, $\vb_\bullet$ & Pesos de entrada, pesos recurrentes y sesgos de cada compuerta. \\
\end{longtable}
}%
}
